\section{Justification}

Industry 4.0 represents a potential business for the future, as evidenced by the substantial growth in investments in related technologies. In 2019 alone, investments in IoT in India reached approximately USD 5 billion \cite{iosr2021iot}. These technologies enable real-time data collection, processing, and transmission, with applications in critical sectors such as energy and water management in developed countries \cite{Water40}.

The growth of the Internet of Things (IoT) is given by the increasing proliferation of connected devices and the adoption of advanced technologies such as 4G, 5G, and Wi-Fi, which rely on OFDM modulation and PSK coding to ensure communication efficiency \cite{10729310}. While these technologies are widely used in daily mobile and office applications, they also face significant challenges due to adverse environmental conditions, including high temperatures, humidity, dust accumulation, and limited maintenance capacity and skilled personnel \cite{8519869}.

A less explored aspect of IoT is its the implementation on high-mobility devices, where equipment operates at variable speeds \cite{10152108}. This variability introduces distortions in the carrier signal, such as ambiguous Doppler effects, which become significant at speeds starting from 25 m/s and distances of 500 m, with a minimum frequency shift of 5 kHz \cite{8022917}. These distortions affect the accuracy of encoded and transmitted data, posing challenges for applications requiring high reliability.

